\label{ch:detector-implementation}

This chapter develops, in full mathematical detail, every model
\texttt{tpeanuts.detector.dayabay} implements: the real multi-reactor
antineutrino flux (\cref{sec:flux-model}), the real order-1/M inverse-beta-decay
cross section (\cref{sec:ibd-cross-section}), the real detector energy
response (\cref{sec:iav-lsnl-theory}), efficiency and exposure
(\cref{sec:efficiency-exposure}), real backgrounds
(\cref{sec:backgrounds}), and the final composed event-rate assembly
(\cref{sec:event-rate-assembly}) with its free-normalization and nuisance
parameters (\cref{sec:nuisances}).

\section{Reactor antineutrino flux}
\label{sec:flux-model}

\modulehead{detector/dayabay/flux.py}{Real per-isotope Huber-Mueller
spectrum, with real non-equilibrium and spent-nuclear-fuel corrections, a
real per-reactor relative power weight, and geometric dilution.}

\begin{theorybox}
\subsection*{Per-isotope spectrum and fission-fraction sum}
The antineutrino spectrum shape at reactor core $r$ is a real
fission-fraction-weighted sum over the four fissile isotopes, with the real
non-equilibrium correction $C_i^{\rm neq}(E)$ folded into each isotope's
tabulated spectrum $S_i(E)$ \emph{before} the sum (published for
$^{235}$U/$^{239}$Pu/$^{241}$Pu only -- no correction exists for
$^{238}$U, which is left uncorrected):
\begin{equation}
  \keyresult{\Sigma(E) = \sum_{i} f_i\, S_i(E)\,\bigl(1+C_i^{\rm neq}(E)\bigr).}
  \label{eq:dayabay-spectrum-shape}
\end{equation}
The non-equilibrium correction accounts for long-lived fission-daughter
beta decays that have not yet reached secular equilibrium with their
production rate; it is largest ($\sim6\%$) just above the IBD threshold and
falls to zero by $\sim4.5$~MeV.

\subsection*{Absolute normalisation and real per-reactor corrections}
The absolute flux at reactor core $r$, before geometric dilution, combines
the real thermal power, the real fission-fraction-weighted mean energy per
fission $\bar e = \sum_i f_i \bar e_i$, a real per-reactor relative
power/output weight $w_r$, and a real per-reactor spent-nuclear-fuel (SNF)
correction $C_r^{\rm SNF}(E)$ (residual antineutrino emission from stored,
previously used fuel assemblies near the core):
\begin{equation}
  \keyresult{\Phi_{{\rm face},r}(E) = \frac{P_{\rm th}}{\bar e}\;
    w_r\;\bigl(1+C_r^{\rm SNF}(E)\bigr)\;\Sigma(E).}
  \label{eq:dayabay-face-flux}
\end{equation}
$w_r$ is derived (not itself a tabulated quantity) from the real weekly
antineutrino-rate history, \cref{sec:weekly-rate-table} below.

\subsection*{Geometric dilution and the multi-reactor sum}
The flux incident on detector $d$ sums the (independently propagated)
contribution of all 6 real cores, each diluted by the real inverse-square
baseline law and multiplied by that reactor's own real oscillation
survival probability:
\begin{equation}
  \keyresult{\Phi_d(E) = \sum_{r=1}^{6}
    \frac{\Phi_{{\rm face},r}(E)}{4\pi L_{r,d}^2}\;
    P_{\bar\nu_e\bar\nu_e}(E,L_{r,d}).}
  \label{eq:dayabay-flux-sum}
\end{equation}
$P_{\bar\nu_e\bar\nu_e}(E,L)$ is the standard vacuum 3-flavour survival
probability, computed once per (reactor, detector) pair at that pair's own
real baseline (\cref{ch:inference-process}); no matter effect is applied
(reactor antineutrinos at these energies and baselines cross negligible
Earth matter). \Cref{eq:dayabay-flux-sum} is evaluated independently for
every one of the 8 real detectors -- each sees its own real 6-baseline
combination.
\end{theorybox}

\subsection*{Real per-reactor relative weight}
\label{sec:weekly-rate-table}

\begin{theorybox}
The official weekly \code{neutrino\_rate} table (\cref{ch:data-release})
gives each real reactor core's own antineutrino emission rate, resolved
week by week across the full 2011--2020 data-taking period, distinguishing
the 6AD/7AD/8AD sub-periods via an \code{n\_det} column. Restricting to the
real 8AD-period weeks and averaging,
\begin{equation}
  \keyresult{w_r = \frac{\left\langle\text{neutrino\_rate}_r\right\rangle_{\rm 8AD}}
    {\frac{1}{6}\sum_{r'} \left\langle\text{neutrino\_rate}_{r'}\right\rangle_{\rm 8AD}},}
  \label{eq:dayabay-relative-weight}
\end{equation}
a 6-reactor average of exactly 1 by construction. This real per-reactor
weight replaces the previous implicit assumption that all 6 cores
contribute identically to \labelcref{eq:dayabay-face-flux} -- the real
measured weights differ from 1 only at the $\sim1$--$2\%$ level (real
Table~\ref{tab:relative-weights-example} in \cref{ch:results}), a small but
genuine, real per-reactor correction.
\end{theorybox}

\begin{implbox}
\begin{itemize}
  \item \pyfunc{reactor\_spectrum\_shape(E\_grid\_MeV)} implements
    \labelcref{eq:dayabay-spectrum-shape} via a cached, corrected copy of
    the raw Huber-Mueller curves (\pyfunc{\_corrected\_hm\_curves}, built
    once with the real non-equilibrium correction interpolated onto each
    isotope's own tabulated energy grid, 0 outside its real defined range).
  \item \pyfunc{reactor\_flux\_at\_face(E\_grid\_MeV, reactor)} implements
    \labelcref{eq:dayabay-face-flux}, with the real SNF correction
    interpolated onto \code{E\_grid\_MeV} (0 outside its own real range).
  \item \pyfunc{flux\_at\_detector(E\_grid\_MeV, baseline\_km, reactor)}
    divides by $4\pi L^2$ (converting \code{baseline\_km} to cm) -- one
    real (reactor, detector) term of \labelcref{eq:dayabay-flux-sum};
    \pyfunc{detector.dayabay.event\_rate.ibd\_event\_rate}
    (\cref{sec:event-rate-assembly}) performs the sum over the 6 real
    reactors explicitly.
  \item \pyfunc{detector.dayabay.parameters.REACTOR\_RELATIVE\_WEIGHT}
    implements \labelcref{eq:dayabay-relative-weight} once, at import time,
    from the real cached weekly-rate table.
\end{itemize}
\end{implbox}

\section{The real order-1/M inverse-beta-decay cross section}
\label{sec:ibd-cross-section}

\modulehead{detector/interaction/inverse\_beta\_decay.py}{Vogel \&
Beacom's threshold-respecting, order-1/M differential IBD cross section
\parencite{VogelBeacom1999}, built by numerical $\cos\theta$ integration
using Daya Bay's own real published coupling constants.}

\begin{theorybox}
\subsection*{Zeroth order (no-recoil) baseline}
At zeroth order in $1/M$ ($M$ the average nucleon mass), the positron
energy is fixed by $E_\nu$ alone,
\begin{equation}
  E_e^{(0)} = E_\nu - \Delta,\qquad \Delta = M_n - M_p,
\end{equation}
and the total cross section is
\begin{equation}
  \sigma^{(0)}(E_\nu) = \sigma_0\,(f^2+3g^2)\,E_e^{(0)} p_e^{(0)},\qquad
  p_e^{(0)} = \sqrt{\bigl(E_e^{(0)}\bigr)^2-m_e^2}.
  \label{eq:ibd-zeroth-order}
\end{equation}
This is Vogel \& Beacom's own Eq.~(10)-(11), and is the formula this
project used for every detector prior to this Daya Bay extension
(\texttt{sigma\_ibd}); it neglects recoil, weak magnetism, and the
resulting positron-angle dependence, and its reaction threshold,
$E_\nu > \bigl((M_n+m_e)^2-M_p^2\bigr)/(2M_p)\approx1.806$~MeV, is exact
only at this order.

\subsection*{Order-1/M differential cross section}
Vogel \& Beacom's own Eq.~(18)-(19) ``high-energy limit'' formula
explicitly neglects the threshold, an effect their own paper states is
``large'' below $\sim30$~MeV -- \emph{every} reactor antineutrino is below
10~MeV, so that formula is not used here. Instead, the full order-1/M
formula (their Eq.~13-15), valid down to threshold, is implemented
directly. At first order the positron energy depends on the
positron-antineutrino scattering angle $\theta$ (lab frame):
\begin{equation}
  \keyresult{E_e^{(1)} = E_e^{(0)}\left[1-\frac{E_\nu}{M}\bigl(1-v_e^{(0)}\cos\theta\bigr)\right]
    - \frac{y^2}{M},\qquad y^2=\frac{\Delta^2-m_e^2}{2},\qquad M=\frac{M_n+M_p}{2},}
  \label{eq:ibd-positron-energy-firstorder}
\end{equation}
with $v_e^{(0)}=p_e^{(0)}/E_e^{(0)}$, and $p_e^{(1)}=\sqrt{(E_e^{(1)})^2-m_e^2}$
(zero if the argument is negative, the exact real threshold cutoff at this
order). The differential cross section becomes
\begin{equation}
  \keyresult{\left(\frac{d\sigma}{d\cos\theta}\right)^{(1)} =
    \frac{\sigma_0}{2}\Bigl[(f^2+3g^2)+(f^2-g^2)\,v_e^{(1)}\cos\theta\Bigr]E_e^{(1)}p_e^{(1)}
    \;-\;\frac{\sigma_0}{2}\,\frac{\Gamma}{M}\,E_e^{(0)}p_e^{(0)},}
  \label{eq:ibd-firstorder-diffxs}
\end{equation}
with $\Gamma$ the weak-magnetism/recoil correction term,
\begin{align}
  \Gamma ={}& 2(f+f_2)g\Bigl[(2E_e^{(0)}+\Delta)(1-v_e^{(0)}\cos\theta)-\tfrac{m_e^2}{E_e^{(0)}}\Bigr] \nonumber\\
  &+ (f^2+g^2)\Bigl[\Delta(1+v_e^{(0)}\cos\theta)+\tfrac{m_e^2}{E_e^{(0)}}\Bigr] \nonumber\\
  &+ (f^2+3g^2)\Bigl[(E_e^{(0)}+\Delta)\bigl(1-\tfrac{\cos\theta}{v_e^{(0)}}\bigr)-\Delta\Bigr] \nonumber\\
  &+ (f^2-g^2)\Bigl[(E_e^{(0)}+\Delta)\bigl(1-\tfrac{\cos\theta}{v_e^{(0)}}\bigr)-\Delta\Bigr]v_e^{(0)}\cos\theta.
  \label{eq:ibd-gamma}
\end{align}

\subsection*{Removing the apparent threshold singularity}
\Cref{eq:ibd-gamma} contains terms $\propto 1/v_e^{(0)}$, apparently
divergent as $E_\nu\to$ threshold ($v_e^{(0)}\to0$); Vogel \& Beacom
themselves note this is exactly cancelled by the explicit $p_e^{(0)}$
prefactor multiplying $\Gamma$ in \labelcref{eq:ibd-firstorder-diffxs}.
Substituting $v_e^{(0)}=p_e^{(0)}/E_e^{(0)}$ throughout \emph{before}
multiplying by $p_e^{(0)}$ makes this cancellation manifest algebraically,
avoiding a numerically risky $0\times\infty$ evaluation at floating-point
precision:
\begin{align}
  p_e^{(0)}\Gamma ={}& 2(f+f_2)g\left[p_e^{(0)}(2E_e^{(0)}+\Delta) -
    \cos\theta\,\frac{\bigl(p_e^{(0)}\bigr)^2}{E_e^{(0)}}(2E_e^{(0)}+\Delta)
    - \frac{p_e^{(0)}m_e^2}{E_e^{(0)}}\right] \nonumber \\
  & + (f^2+g^2)\left[\Delta p_e^{(0)} + \cos\theta\,\Delta\,\frac{\bigl(p_e^{(0)}\bigr)^2}{E_e^{(0)}}
    + \frac{p_e^{(0)}m_e^2}{E_e^{(0)}}\right] \nonumber \\
  & + (f^2+3g^2)\Bigl[p_e^{(0)}E_e^{(0)}-\cos\theta\,E_e^{(0)}(E_e^{(0)}+\Delta)\Bigr] \nonumber \\
  & + (f^2-g^2)\,p_e^{(0)}\cos\theta\Bigl[p_e^{(0)}-\cos\theta(E_e^{(0)}+\Delta)\Bigr],
  \label{eq:ibd-gamma-clean}
\end{align}
free of any $1/v_e^{(0)}$ or $1/p_e^{(0)}$ term -- every division is by
$E_e^{(0)}$, which is bounded below by $m_e>0$ over the whole real IBD
kinematic range.

\subsection*{Bare normalisation $\sigma_0$}
The bare $\sigma_0$ (Vogel \& Beacom's Eq.~9) does not depend on
$f/g/f_2$; it is computed here directly rather than back-derived from the
zeroth-order coefficient (which was calibrated using Vogel \& Beacom's own,
slightly older $g=1.26$ value, not Daya Bay's real published
$g=1.2701$):
\begin{equation}
  \keyresult{\sigma_0 = \frac{G_F^2\cos^2\theta_C}{\pi}\bigl(1+\Delta_{\rm inner}^R\bigr)
    \times (\hbar c)^2,}
  \label{eq:ibd-sigma0-firstprinciples}
\end{equation}
with $\cos\theta_C=0.974$ the Cabibbo angle and $\Delta_{\rm inner}^R\approx
0.024$ the energy-independent inner radiative correction, both from
\cite{VogelBeacom1999}, and $(\hbar c)^2$ converting from natural units to
$\mathrm{cm^2/MeV^4}$. This reproduces Vogel \& Beacom's own quoted
$0.0952\times10^{-42}\ \mathrm{cm^2/MeV^2}$ coefficient to $\sim1\%$ when
recombined with their own $f=1$, $g=1.26$ -- the residual difference is
input-value vintage (rounding of $\Delta_{\rm inner}^R$/$\cos\theta_C$),
not a formula error.

\subsection*{Total cross section and the prompt-energy differential}
The total cross section is obtained by numerical integration over
$\cos\theta$,
\begin{equation}
  \sigma_{\rm tot}(E_\nu) = \int_{-1}^{1} d\cos\theta \left(\frac{d\sigma}{d\cos\theta}\right)^{(1)},
\end{equation}
the same method Vogel \& Beacom themselves used for their own published
curves. To build the differential cross section in the observable prompt
energy $d\sigma/dT$, each $\cos\theta$-quadrature node's cross-section
contribution is redistributed onto the prompt-energy grid at its own
order-1/M prompt energy,
\begin{equation}
  E_{\rm prompt} = E_e^{(1)} + m_e,
\end{equation}
via linear-interpolation mass conservation (identical technique to
\labelcref{eq:iav-warp-weights} below), then converted from a
per-quadrature-node mass to a density by dividing by the prompt-energy grid
spacing.
\end{theorybox}

\begin{implbox}
\begin{itemize}
  \item \pyfunc{sigma\_ibd\_diff\_cos\_precise(E\_nu\_MeV, cos\_theta, f=1.0,
    g=1.2701, f2=3.706)} implements
    \labelcref{eq:ibd-positron-energy-firstorder,eq:ibd-firstorder-diffxs,eq:ibd-gamma-clean}
    directly, defaulting to Daya Bay's own real coupling constants
    (\code{ibd\_constants.csv}, \cref{ch:data-release}); returns
    $(d\sigma/d\cos\theta,\,E_e^{(1)})$.
  \item \pyfunc{ibd\_cross\_section\_grid\_precise(E\_nu\_grid\_MeV,
    T\_grid\_MeV, n\_cos=41)} performs the $\cos\theta$-quadrature-and-scatter
    construction above on a uniform 41-point trapezoidal grid, returning
    $d\sigma/dT$ on the outer-product grid
    \code{(E\_nu\_grid\_MeV, T\_grid\_MeV)}, shape \code{(n\_E,n\_T)},
    $\mathrm{cm^2/MeV}$.
  \item \pyfunc{\_sigma0\_bare\_cm2\_per\_mev4()} implements
    \labelcref{eq:ibd-sigma0-firstprinciples}, cached at module import.
  \item \pyfunc{sigma\_ibd(E\_nu\_MeV)}/\pyfunc{ibd\_cross\_section\_grid(...)}:
    the original zeroth-order formula, \labelcref{eq:ibd-zeroth-order},
    retained unchanged (used by any detector that has not opted into the
    precise cross section).
  \item \pyfunc{detector.common.response.scatter\_add\_linear(grid,
    positions, values)}: the shared mass-conserving linear-interpolation
    scatter utility used both here (redistributing cross-section mass
    across $\cos\theta$ nodes) and by the LSNL warp matrix
    (\cref{sec:iav-lsnl-theory}).
\end{itemize}
\end{implbox}

\begin{notebox}
The real, measured effect of this upgrade: the order-1/M cross section is a
few percent \emph{below} the zeroth-order one, growing with $E_\nu$
(consistent with Vogel \& Beacom's own stated $\mathcal O(E_\nu/M)$
scaling, with an $\mathcal O(-7)$ numerical coefficient) -- validated
independently in this project's own test suite
(\code{detector/dayabay/test/test2\_response\_and\_fit.py}).
\end{notebox}

\section{Detector energy response: IAV, LSNL, and Gaussian resolution}
\label{sec:iav-lsnl-theory}

\modulehead{detector/dayabay/response.py}{The full real response chain,
applied in the official Daya Bay Collaboration analysis pipeline's own
order: IAV redistribution, then LSNL nonlinearity, then Gaussian
photostatistics resolution.}

\begin{theorybox}
\subsection*{Real response order}
The official Daya Bay dagflow/GNA analysis pipeline's own node names
(confirmed directly from the Collaboration's own CHEP~2026 conference
presentation on that framework) read, cumulatively,
\begin{equation}
  \text{raw} \;\longrightarrow\; \text{stages.iav} \;\longrightarrow\;
  \text{stages.evis (IAV+LSNL)} \;\longrightarrow\; \text{stages.erec (+resolution)},
\end{equation}
i.e.\ \keyresult{IAV is applied to the true deposited energy \emph{before}
LSNL}, which is then followed by the Gaussian resolution smearing. This
also matches the physical picture: energy lost into the inert acrylic
vessel material near its wall (the IAV effect) never produces
scintillation light, so it cannot be subject to a light-yield nonlinearity
curve (LSNL) -- only the light-producing, IAV-redistributed energy is.

\subsection*{IAV redistribution}
The real IAV correction is a $240\times240$ energy-redistribution matrix
$R^{\rm IAV}$ on a $0.05$~MeV, $0$--$12$~MeV grid (matching this project's
own \code{T\_GRID\_MEV} exactly), normalised so that each column sums to
1 -- a genuine discrete probability-mass (not density) transfer matrix,
\begin{equation}
  R^{\rm IAV}_{ij} = P(\text{reconstructed in bin } i \mid \text{true energy in bin } j),
  \qquad \sum_i R^{\rm IAV}_{ij} = 1.
\end{equation}

\subsection*{LSNL nonlinearity}
The real LSNL correction is a multiplicative energy-scale warp,
$f(E)=E_{\rm reconstructed}/E_{\rm true\ deposited}$, mapping true energy
$T$ to $T_{\rm nl}=T\cdot f(T)$. Represented as a mass-conserving matrix on
the same grid, column $j$ (true energy $T_j$) is split by linear
interpolation between the two grid points bracketing $T_{\rm nl}(T_j)$:
\begin{equation}
  \keyresult{R^{\rm LSNL}_{ij} = \begin{cases}
    1-w_j, & i = \lfloor T_{\rm nl}(T_j)\rfloor_{\rm grid}, \\
    w_j, & i = \lceil T_{\rm nl}(T_j)\rceil_{\rm grid}, \\
    0, & \text{otherwise,}
  \end{cases}}
  \qquad w_j = \frac{T_{\rm nl}(T_j)-T_{\lfloor\cdot\rfloor}}{T_{\lceil\cdot\rceil}-T_{\lfloor\cdot\rfloor}}
  \label{eq:iav-warp-weights}
\end{equation}
(clamped to $[0,1]$, with all mass dumped at the nearest edge for a
position outside the grid range) -- the identical construction used for
the IBD cross section's $\cos\theta\to T$ redistribution
(\cref{sec:ibd-cross-section}). By construction $\sum_i R^{\rm LSNL}_{ij}=1$
for every column $j$.

\subsection*{Official LSNL pull-curve model}
The real LSNL curve used is not fixed to its nominal shape alone: the
official systematic model is a linear combination of the nominal curve
$f_0$ and 4 real published pull curves $f_k$,
\begin{equation}
  \keyresult{f(E) = f_0(E) + \sum_{k=1}^{4} a_k\bigl(f_k(E)-f_0(E)\bigr),}
  \label{eq:lsnl-pull-model}
\end{equation}
with each $a_k$ a real free nuisance carrying an official independent,
zero-mean, unit-sigma Gaussian prior (\code{parameters/detector\_lsnl.yaml}).
At $a_k=0$ for all $k$, \labelcref{eq:lsnl-pull-model} reduces to the
nominal curve; \cref{ch:inference-process} documents how $a_k$ enters the
fit as a real, Gaussian-constrained nuisance parameter.

\subsection*{Gaussian photostatistics resolution}
The real 3-term resolution formula,
\begin{equation}
  \keyresult{\frac{\sigma(E)}{E} = \sqrt{a^2+\frac{b^2}{E}+\frac{c^2}{E^2}}},
  \qquad a=0.016,\ b=0.081\ \sqrt{\rm MeV},\ c=0.026\ {\rm MeV},
  \label{eq:eres-formula}
\end{equation}
with terms interpretable as spatial/temporal non-uniformity ($a$), photon
(photoelectron) counting statistics ($b$), and PMT dark noise/electronics
($c$), builds a genuine probability \emph{density} in $T'$ (not a discrete
mass matrix, unlike IAV/LSNL): $R^{\rm Gauss}(T'|T) =
\mathcal N(T';T,\sigma(T))$.

\subsection*{Composition}
All three steps are real, independent, official detector-response
ingredients (the data release ships them as separate real files precisely
because they are physically distinct effects); they are composed as
sequential linear operators. Because the folding machinery
(\cref{sec:event-rate-assembly}) integrates $R(T'|T)$ against $T$ via
trapezoidal quadrature, every step is first built as a discrete,
mass-conserving transfer matrix on the shared 0.05~MeV grid, and only the
final \emph{composed} matrix is converted from probability mass to the
density the folding integral expects (dividing by the grid spacing $\delta
T$):
\begin{equation}
  \keyresult{R(T'|T) = \frac{1}{\delta T}\,
    R^{\rm Gauss}_{\rm mass}\;R^{\rm LSNL}\;R^{\rm IAV},}
  \label{eq:response-composition}
\end{equation}
where $R^{\rm Gauss}_{\rm mass} = \delta T\cdot R^{\rm Gauss}$ (density
converted to an approximate mass for the composition, valid since $\delta
T=0.05$~MeV is fine compared to $\sigma(E)$ over the whole relevant
energy range). Reading \labelcref{eq:response-composition} as a matrix
product applied to a column vector $v$ representing a true-energy
spectrum, $R(T'|T)\,v = R^{\rm Gauss}_{\rm mass}\bigl(R^{\rm LSNL}(R^{\rm
IAV}v)\bigr)/\delta T$ -- $v$ is transformed by $R^{\rm IAV}$ first, then
$R^{\rm LSNL}$, then $R^{\rm Gauss}_{\rm mass}$, exactly the real physical
order stated above.
\end{theorybox}

\begin{implbox}
\begin{itemize}
  \item \pyfunc{sigma\_MeV(T\_grid\_MeV, a, b, c)} implements
    \labelcref{eq:eres-formula}.
  \item \pyfunc{\_iav\_mass\_matrix(T\_grid\_MeV)} embeds the real
    $240\times240$ \pyfunc{IAV\_MATRIX} into the project's own (possibly
    larger) grid, with an identity fallback for any extra grid point beyond
    the real IAV table's own range (physically negligible: the real IBD
    prompt-energy spectrum is essentially zero there).
  \item \pyfunc{\_lsnl\_curves\_on\_grid(T\_grid\_MeV)} interpolates the
    real nominal/pull curves onto the shared grid once
    (\code{@torch.no\_grad()}, fixed real data).
  \item \pyfunc{\_lsnl\_warp\_matrix(T\_grid\_MeV, lsnl\_pulls=None)}
    implements \labelcref{eq:lsnl-pull-model,eq:iav-warp-weights};
    differentiable w.r.t.\ \code{lsnl\_pulls} (only the linear combination
    and the resulting warp weights carry gradient -- the underlying curve
    interpolation is fixed real data).
  \item \pyfunc{response\_matrix(T\_grid\_MeV, Tprime\_grid\_MeV, a, b, c,
    lsnl\_pulls=None)} implements \labelcref{eq:response-composition}.
\end{itemize}
\end{implbox}

\begin{notebox}
Every component above uses the real official nominal/central-value
curves/matrix; the LSNL pull-curve nuisances $a_k$ can be fit
(\cref{ch:inference-process}), but the IAV matrix's own official
off-diagonal-scale systematic (\code{detector\_iav\_offdiag\_scale.yaml},
a single real nuisance with a $4\%$ prior) is not implemented. The three
response steps are composed assuming they act as independent, sequential
linear operators -- a standard simplification when a full joint response
covariance is not being reproduced.
\end{notebox}

\section{Efficiency and exposure}
\label{sec:efficiency-exposure}

\begin{theorybox}
The real IBD-selection efficiency $\varepsilon\approx0.8025$
(Gd-capture fraction $\times$ analysis-cut efficiencies, consistent with
the real published individual cut efficiencies -- Gd-capture fraction
$\approx85.5\%$, prompt-energy cut $\approx99.81\%$, capture-time cut
$\approx98.70\%$, multiplicity cut $\approx97.5\%$, flasher cut
$\approx99.98\%$, whose product is $\approx81.6\%$, close to the quoted
$80.25\%$) is applied as a flat multiplicative factor, on top of and
separate from the real per-detector, per-day effective livetime
(\code{eff\_livetime} $=$ \code{livetime}$\times$\code{eff}, the latter
itself the real combined muon-veto and multiplicity-cut daily efficiency),
summed over the real 8AD-period days to give each detector's total real
exposure in seconds.
\end{theorybox}

\begin{implbox}
\pyfunc{DETECTOR\_EFFICIENCY} (a flat scalar) multiplies the folded
reconstructed spectrum after response and before binning
(\cref{sec:event-rate-assembly}); \pyfunc{load\_exposure()} sums
\code{eff\_livetime} over real 8AD-period days per detector.
\end{implbox}

\section{Real backgrounds}
\label{sec:backgrounds}

\begin{theorybox}
Five real background categories contribute to the observed spectrum,
each with its own real published per-detector rate (events/day) and real
normalised shape histogram (integrating to 1 on the same 0.05~MeV binning
as the signal):
\begin{itemize}
  \item \textbf{Accidentals} -- uncorrelated prompt-like/delayed-like
    triggers randomly coincident in time; the real best-measured category
    (fractional rate uncertainty $\sim0.1\%$).
  \item \textbf{$^9$Li/$^8$He} -- cosmogenic isotopes produced by
    through-going muons, beta-decaying with a delayed-neutron-emission
    branch that mimics the IBD signature.
  \item \textbf{Fast neutrons} -- energetic neutrons from muon spallation
    outside the water shield, which can deposit a prompt-like recoil energy
    followed by a real neutron capture.
  \item \textbf{AmC} -- correlated background from the AmC neutron
    calibration sources themselves.
  \item \textbf{$^{13}$C$(\alpha,n)^{16}$O} -- alpha decays of natural
    radioactivity in the scintillator inducing a real $(\alpha,n)$ reaction
    on carbon.
\end{itemize}
The predicted background count per real analysis bin is
\begin{equation}
  N_k^{\rm bkg} = \sum_{\rm category} ({\rm rate}_{\rm category}\times{\rm exposure\_days})
    \times{\rm shape}_{{\rm category},k},
\end{equation}
with the fine-binned (0.05~MeV) product rebinned onto the real
non-uniform analysis edges by a plain per-bin sum (every real analysis-bin
edge lands exactly on a fine-bin boundary, so no fine bin straddles two
coarse bins).
\end{theorybox}

\begin{implbox}
\pyfunc{real\_background\_counts(detector, exposure\_seconds,
bin\_edges\_MeV, category\_scale=None)} implements the sum above;
\pyfunc{rebin\_discrete\_counts} performs the fine-to-coarse rebinning.
An optional \code{category\_scale} dict multiplies each category's
contribution by a real per-category nuisance (\cref{sec:nuisances}).
\end{implbox}

\section{Composed event-rate assembly}
\label{sec:event-rate-assembly}

\modulehead{detector/dayabay/event\_rate.py}{The one place this package's
real 6-reactor geometry, real flux, real order-1/M cross section, real
response, and real background come together, per detector.}

\begin{theorybox}
For one real detector $d$, the predicted signal (pre-response) spectrum
folds the multi-reactor flux sum \labelcref{eq:dayabay-flux-sum} with the
order-1/M differential cross section \labelcref{eq:ibd-firstorder-diffxs},
\begin{equation}
  \frac{dR}{dT}(T) = N_p^{(d)}\int dE_\nu\;\Phi_d(E_\nu)\,\frac{d\sigma}{dT}(E_\nu,T),
\end{equation}
with $N_p^{(d)}$ the real target free-proton count for detector $d$; the
response matrix \labelcref{eq:response-composition} is then applied,
\begin{equation}
  \frac{dR}{dT'}(T') = \int dT\; R(T'|T)\,\frac{dR}{dT}(T),
\end{equation}
followed by the flat efficiency factor and integration into the real
non-uniform analysis bins,
\begin{equation}
  \keyresult{N_k(\theta) = \nu_{\rm norm}\cdot{\rm exposure}\cdot\varepsilon\int_{{\rm bin}\,k}
    dT'\;\frac{dR}{dT'}(T';\theta) \;+\; N_k^{\rm bkg},}
  \label{eq:dayabay-full-prediction}
\end{equation}
the full, real, per-bin, per-detector prediction, with $\nu_{\rm norm}$
the free global-normalization nuisance (\cref{sec:nuisances}) multiplying
the signal only (the real background is separately measured and not
rescaled by it).
\end{theorybox}

\begin{implbox}
\begin{itemize}
  \item \pyfunc{ibd\_event\_rate(detector, p\_ee\_per\_reactor, ...,
    signal\_scale=1.0, background\_category\_scale=None, lsnl\_pulls=None)}
    implements \labelcref{eq:dayabay-full-prediction} end to end, calling
    \pyfunc{detector.common.event\_rate.predicted\_counts} with the real
    cross section, real response, real target, real exposure, and real
    background plugged in.
  \item \code{p\_ee\_per\_reactor}: \code{\{reactor: P\_ee(E)\}}, one
    tensor per real reactor, already oscillated at that reactor's own
    real baseline -- supplied by
    \pyclass{detector.dayabay.inference\_model.DayaBayDetectorModel}
    (\cref{ch:inference-process}), which itself calls
    \texttt{tpeanuts.medium.vacuum.oscillation\_model
    .VacuumOscillationModel.predict\_pee} once per reactor.
\end{itemize}
\end{implbox}

\section{Free normalization and nuisance parameters}
\label{sec:nuisances}

\begin{theorybox}
\subsection*{Global normalization}
Daya Bay's own official combined analysis leaves a free overall
normalization nuisance in its fit (\code{parameters/detector\_normalization.yaml},
real nominal value 1.0, labelled ``Global IBD rate normalization''): even
Daya Bay's own simulation chain is not asserted to self-normalize to the
observed rate without fitting it. \Cref{eq:dayabay-full-prediction}'s
$\nu_{\rm norm}$ is this same real, official parameter, not an ad hoc
rescaling invented for this project -- see \cref{ch:results} for why it is
needed (a real, documented $\sim19\%$ absolute-normalization gap, audited
and explained in \cref{ch:results}, not attributable to any single unit or
loading error).

\subsection*{Real background-rate nuisances}
Five further real nuisances, one Gaussian-constrained multiplicative scale
$c_k$ per background category (nominal 1), with real prior width
\begin{equation}
  \sigma_k = \frac{\sqrt{\sum_d ({\rm uncertainty}_{k,d})^2}}{\sum_d {\rm rate}_{k,d}},
  \label{eq:background-sigma}
\end{equation}
a rate-weighted quadrature combination of the real per-detector
(rate, uncertainty) pairs into one correlated scale per category -- a
simplification of the official per-detector/per-hall correlation structure
(accidentals genuinely uncorrelated between detectors; AmC correlated
across all detectors; $^9$Li/$^8$He and fast neutrons correlated within a
hall), not a per-detector free parameter.

\subsection*{Real LSNL pull-curve nuisances}
The 4 real $a_k$ of \labelcref{eq:lsnl-pull-model}, official prior
independent zero-mean unit-sigma Gaussian.
\end{theorybox}

\begin{implbox}
\begin{itemize}
  \item \pyclass{DayaBayDetectorModel}(\code{normalization\_free},
    \code{background\_free}, \code{lsnl\_free}): each flag appends the
    corresponding real nuisance name(s) to \code{.free} and slices them
    from \code{theta} in \code{.predict} (\cref{ch:inference-process}).
  \item \pyfunc{detector.dayabay.parameters.BACKGROUND\_CATEGORY\_SIGMA}
    implements \labelcref{eq:background-sigma} once, at import time.
  \item \pyfunc{tpeanuts.inference.likelihood.gaussian\_prior\_penalty(values,
    mean, sigma)}: the generic, reusable $-2\ln L$ Gaussian-prior penalty
    added to the fit's likelihood for any of these nuisances
    (\cref{ch:inference-process}).
\end{itemize}
\end{implbox}
